\documentclass[11pt,a4paper]{article}
% preamble.tex — estilo común de todos los apuntes Froth.
% Cada apunte hace % preamble.tex — estilo común de todos los apuntes Froth.
% Cada apunte hace % preamble.tex — estilo común de todos los apuntes Froth.
% Cada apunte hace \input{preamble} justo después de \documentclass.
\usepackage[margin=2.4cm]{geometry}
\usepackage{xcolor}
\definecolor{litblue}{RGB}{31,79,121}
\definecolor{litgreen}{RGB}{27,94,32}
\definecolor{litorange}{RGB}{230,81,0}
\usepackage{parskip}
\usepackage{titlesec}
\titleformat{\section}{\Large\bfseries\color{litblue}}{\thesection}{0.6em}{}
\titleformat{\subsection}{\large\bfseries\color{litblue}}{\thesubsection}{0.6em}{}
\usepackage{tcolorbox}
\usepackage{fancyhdr}
\pagestyle{fancy}
\fancyhf{}
\fancyhead[L]{\small\color{litblue}Froth — Apuntes de ML}
\fancyhead[R]{\small\thepage}
\renewcommand{\headrulewidth}{0.4pt}
\usepackage[colorlinks=true,linkcolor=litblue,urlcolor=litblue]{hyperref}

% Cabecera de cada apunte: \cabecera{numero}{titulo}
\newcommand{\cabecera}[2]{%
  \begin{tcolorbox}[colback=litblue,coltext=white,colframe=litblue,arc=2mm]
    {\Large\bfseries Apunte #1 — #2}\\[3pt]
    {\small Proyecto Froth · aprender Machine Learning construyendo}
  \end{tcolorbox}\vspace{0.8em}}

% Cajas temáticas
\newtcolorbox{concepto}[1]{colback=blue!4,colframe=litblue,title={Concepto: #1},arc=1.5mm}
\newtcolorbox{practica}{colback=green!5,colframe=litgreen,title={Pruébalo tú (teclado en mano)},arc=1.5mm}
\newtcolorbox{ojo}{colback=orange!8,colframe=litorange,title={Ojo / error típico},arc=1.5mm}
\newtcolorbox{resumen}{colback=gray!8,colframe=gray!60,title={En una frase},arc=1.5mm}

\newcommand{\codigo}[1]{\texttt{#1}}
 justo después de \documentclass.
\usepackage[margin=2.4cm]{geometry}
\usepackage{xcolor}
\definecolor{litblue}{RGB}{31,79,121}
\definecolor{litgreen}{RGB}{27,94,32}
\definecolor{litorange}{RGB}{230,81,0}
\usepackage{parskip}
\usepackage{titlesec}
\titleformat{\section}{\Large\bfseries\color{litblue}}{\thesection}{0.6em}{}
\titleformat{\subsection}{\large\bfseries\color{litblue}}{\thesubsection}{0.6em}{}
\usepackage{tcolorbox}
\usepackage{fancyhdr}
\pagestyle{fancy}
\fancyhf{}
\fancyhead[L]{\small\color{litblue}Froth — Apuntes de ML}
\fancyhead[R]{\small\thepage}
\renewcommand{\headrulewidth}{0.4pt}
\usepackage[colorlinks=true,linkcolor=litblue,urlcolor=litblue]{hyperref}

% Cabecera de cada apunte: \cabecera{numero}{titulo}
\newcommand{\cabecera}[2]{%
  \begin{tcolorbox}[colback=litblue,coltext=white,colframe=litblue,arc=2mm]
    {\Large\bfseries Apunte #1 — #2}\\[3pt]
    {\small Proyecto Froth · aprender Machine Learning construyendo}
  \end{tcolorbox}\vspace{0.8em}}

% Cajas temáticas
\newtcolorbox{concepto}[1]{colback=blue!4,colframe=litblue,title={Concepto: #1},arc=1.5mm}
\newtcolorbox{practica}{colback=green!5,colframe=litgreen,title={Pruébalo tú (teclado en mano)},arc=1.5mm}
\newtcolorbox{ojo}{colback=orange!8,colframe=litorange,title={Ojo / error típico},arc=1.5mm}
\newtcolorbox{resumen}{colback=gray!8,colframe=gray!60,title={En una frase},arc=1.5mm}

\newcommand{\codigo}[1]{\texttt{#1}}
 justo después de \documentclass.
\usepackage[margin=2.4cm]{geometry}
\usepackage{xcolor}
\definecolor{litblue}{RGB}{31,79,121}
\definecolor{litgreen}{RGB}{27,94,32}
\definecolor{litorange}{RGB}{230,81,0}
\usepackage{parskip}
\usepackage{titlesec}
\titleformat{\section}{\Large\bfseries\color{litblue}}{\thesection}{0.6em}{}
\titleformat{\subsection}{\large\bfseries\color{litblue}}{\thesubsection}{0.6em}{}
\usepackage{tcolorbox}
\usepackage{fancyhdr}
\pagestyle{fancy}
\fancyhf{}
\fancyhead[L]{\small\color{litblue}Froth — Apuntes de ML}
\fancyhead[R]{\small\thepage}
\renewcommand{\headrulewidth}{0.4pt}
\usepackage[colorlinks=true,linkcolor=litblue,urlcolor=litblue]{hyperref}

% Cabecera de cada apunte: \cabecera{numero}{titulo}
\newcommand{\cabecera}[2]{%
  \begin{tcolorbox}[colback=litblue,coltext=white,colframe=litblue,arc=2mm]
    {\Large\bfseries Apunte #1 — #2}\\[3pt]
    {\small Proyecto Froth · aprender Machine Learning construyendo}
  \end{tcolorbox}\vspace{0.8em}}

% Cajas temáticas
\newtcolorbox{concepto}[1]{colback=blue!4,colframe=litblue,title={Concepto: #1},arc=1.5mm}
\newtcolorbox{practica}{colback=green!5,colframe=litgreen,title={Pruébalo tú (teclado en mano)},arc=1.5mm}
\newtcolorbox{ojo}{colback=orange!8,colframe=litorange,title={Ojo / error típico},arc=1.5mm}
\newtcolorbox{resumen}{colback=gray!8,colframe=gray!60,title={En una frase},arc=1.5mm}

\newcommand{\codigo}[1]{\texttt{#1}}

\begin{document}
\cabecera{20}{La red de similitud: K vecinos, umbral, y tus hubs y puentes reales}

\section{Cómo se construyó la red}

Cada paper se conecta con sus \textbf{K papers más parecidos} (similitud coseno calculada en
el espacio \textbf{original de 768D} — no en el mapa 2D, que ya sabemos que distorsiona), y
solo sobreviven las conexiones con similitud $\geq$ un \textbf{umbral}. Es el llamado
\emph{grafo kNN} (k nearest neighbors).

\begin{concepto}{las perillas K y umbral: nadie las ``sabe'', se prueban}
Preguntaste: ¿y quién dice que 5 vecinos es lo ideal? Respuesta honesta: \textbf{nadie} — es
una perilla, como la granularidad. La forma seria de elegirla es probar valores y mirar la
estructura resultante:

\begin{center}
\begin{tabular}{rrrrr}
\textbf{K} & \textbf{líneas} & \textbf{conex./paper} & \textbf{piezas} & \textbf{\% entre clusters} \\
\hline
3 & 281 & 4.7 & 1 & 17\% \\
\textbf{5} & \textbf{457} & \textbf{7.6} & \textbf{1} & \textbf{21\%} \\
8 & 702 & 11.7 & 1 & 24\% \\
15 & 1213 & 20.2 & 1 & 30\% \\
\end{tabular}
\end{center}

K alto revela más cruces entre temas (tus primeros vecinos son de tu familia; los lejanos
cruzan fronteras) pero con 20 conexiones por paper el dibujo es ilegible (``bola de pelo'').
Elegido: \textbf{K=5, umbral 0.75} — legible, con puentes visibles y la red en una sola
pieza. Ambos viven en \codigo{config.py} (\codigo{GRAPH\_KNN},
\codigo{GRAPH\_SIM\_THRESHOLD}): cambiar de opinión cuesta una línea.
\end{concepto}

\section{Medir hubs y puentes (ya no a ojo)}

\begin{concepto}{grado y betweenness}
\begin{itemize}
  \item \textbf{Hub} $=$ nodo con muchas conexiones $\rightarrow$ se mide con el
        \textbf{grado} (contar sus aristas).
  \item \textbf{Puente} $=$ nodo por el que pasan los caminos entre comunidades $\rightarrow$
        se mide con la \textbf{betweenness centrality}: cuántas veces un nodo está en el
        camino más corto entre otros dos. Un paper con \emph{pocas} conexiones puede ser un
        gran puente si esas pocas unen mundos.
\end{itemize}
\end{concepto}

\section{Los resultados con TUS papers}

\textbf{Hubs} (las lecturas obligatorias del corpus):
\begin{enumerate}
  \item \emph{The Magmatic--Hydrothermal Transition in Lithium Pegmatites} (24 conexiones)
  \item \emph{Electrostatically Controlled Enrichment of Lepidolite via Flotation} (15)
  \item \emph{Mineral Chemistry of Two-Mica Granite Rare Metals} (15)
\end{enumerate}

\textbf{Puentes} (donde vive la novedad) — con dos hallazgos gordos:
\begin{enumerate}
  \item Entre los 5 puentes top está \emph{``Characterization and Liberation Study of the
        \textbf{Beauvoir Granite}''}: el paper fundacional de TU yacimiento es un puente
        entre comunidades. La red confirma \textbf{con números} que el territorio de tu
        tesis (geología $\leftrightarrow$ procesamiento en Beauvoir) es terreno-frontera:
        argumento de originalidad medible para tu introducción.
  \item 3 de los 5 puentes top estaban etiquetados como \textbf{ruido} por HDBSCAN. La
        lección del Apunte 14 (``ruido $\neq$ basura; son generalistas entre nubes'') se
        cumplió cuantitativamente: \emph{el ruido de ayer son los puentes de hoy}.
\end{enumerate}

\section{Dónde vive en el código}

\codigo{froth/graph.py}: \codigo{build\_similarity\_graph()} (la red, con atributos por
nodo), \codigo{find\_hubs()} (grado), \codigo{find\_bridges()} (betweenness),
\codigo{save\_graph()/load\_graph()} (formato GraphML en \codigo{2\_Datos/processed/}).
Se corre con \codigo{python -m froth.graph}.

\begin{resumen}
Red kNN: cada paper con sus K=5 más parecidos (coseno en 768D), umbral 0.75 — perillas
elegidas probando, guardadas en config. Grado encuentra hubs; betweenness encuentra puentes.
Tu corpus confirmó: el paper de Beauvoir es un puente, y el ``ruido'' escondía a los puentes.
\end{resumen}

\end{document}
